% ============================================================================
% MUSTERLÖSUNG: Reibungskräfte (komplett - Doppelstunde)
% Physik Klasse 7 - Kräfte
% Stand: 05.02.2026
% ============================================================================

\documentclass[11pt,a4paper]{article}

\usepackage[utf8]{inputenc}
\usepackage[T1]{fontenc}
\usepackage[ngerman]{babel}
\usepackage{geometry}
\usepackage{helvet}
\usepackage{enumitem}
\usepackage{fancyhdr}
\usepackage{xcolor}
\usepackage{tcolorbox}
\usepackage{amssymb}
\usepackage{siunitx}
\usepackage{array}
\usepackage{tabularx}
\usepackage{textcomp}
\usepackage[hidelinks]{hyperref}

\sisetup{locale=DE, per-mode=power, inter-unit-product=\ensuremath{\cdot}}
\renewcommand{\familydefault}{\sfdefault}

\geometry{left=1.4cm, right=1.4cm, top=1.3cm, bottom=1.0cm}

% Fachfarbe Physik
\definecolor{physik}{HTML}{2c5aa0}
\colorlet{fachfarbe}{physik}
\definecolor{protokollrand}{HTML}{2a7a4b}
\definecolor{merkerand}{HTML}{b35c00}
\definecolor{loesung}{HTML}{c62828}

% Kopfzeile
\pagestyle{fancy}
\fancyhf{}
\fancyhead[L]{\small\textbf{Physik Klasse 7} | Reibungskräfte}
\fancyhead[R]{\small\textcolor{loesung}{\textbf{MUSTERLÖSUNG}}}
\renewcommand{\headrulewidth}{0.4pt}
\setlength{\headheight}{14pt}

% Boxen
\tcbuselibrary{skins}

\newtcolorbox{merkkasten}[1][]{
    colback=white, colframe=fachfarbe,
    coltitle=black, colbacktitle=white,
    fonttitle=\bfseries\small,
    boxrule=1pt, arc=0pt,
    left=6pt, right=6pt, top=4pt, bottom=4pt,
    #1
}

\newtcolorbox{protokollbox}[1][]{
    colback=white, colframe=protokollrand,
    coltitle=black, colbacktitle=white,
    fonttitle=\bfseries\small,
    boxrule=1pt, arc=0pt,
    left=6pt, right=6pt, top=4pt, bottom=4pt,
    #1
}

\newtcolorbox{merkebox}[1][]{
    colback=white, colframe=merkerand,
    coltitle=black, colbacktitle=white,
    fonttitle=\bfseries\small,
    boxrule=1pt, arc=0pt,
    left=6pt, right=6pt, top=4pt, bottom=4pt,
    #1
}

\newcommand{\aufgabe}[1]{\noindent\textbf{\textcolor{fachfarbe}{#1}}}
\newcommand{\lsg}[1]{\textcolor{loesung}{\textbf{#1}}}
\newcommand{\stern}[1]{\textcolor{merkerand}{\textbf{#1}}}

\setlength{\parindent}{0pt}
\setlength{\parskip}{0pt}

\begin{document}

% ============================================================================
% TITEL
% ============================================================================
{\Large\bfseries Reibungskräfte -- \textcolor{loesung}{Musterlösung}}\\[0.2em]
\hrule
\vspace{0.5em}

% ============================================================================
% B1: BEOBACHTUNG HAFT/GLEIT
% ============================================================================
\aufgabe{B1: Beobachtung -- Klotz anschieben}

\vspace{0.2em}
Um den Klotz in Bewegung zu setzen, muss die Kraft erst \lsg{einen Schwellenwert überschreiten / die Haftreibung überwinden}.

Sobald der Klotz gleitet, \lsg{sinkt / verringert sich} die benötigte Kraft.

\vspace{0.2em}

% ============================================================================
% H1: HYPOTHESE
% ============================================================================
\aufgabe{H1: Hypothese -- Meine Vermutung}

\vspace{0.2em}
\lsg{$\boxtimes$ Haftreibung $>$ Gleitreibung $>$ Rollreibung} {\small (richtige Antwort)}

\vspace{0.3em}

% ============================================================================
% T1: MESSWERTE
% ============================================================================
\begin{protokollbox}[title=T1: Messwerte -- Drei Reibungsarten]
\small
\renewcommand{\arraystretch}{1.3}
\begin{tabular}{|p{4cm}|p{3cm}|p{5cm}|}
\hline
\textbf{Reibungsart} & \textbf{Messwert} & \textbf{Beschreibung} \\
\hline
Haftreibung (Losbrechkraft) & \lsg{8} N & Kraft zum Lösen \\
\hline
Gleitreibung & \lsg{5} N & Kraft während Bewegung \\
\hline
Rollreibung & \lsg{1} N & Kraft beim Rollen \\
\hline
\end{tabular}
\end{protokollbox}

\vspace{0.2em}

% ============================================================================
% B2/K2: ERKENNTNIS
% ============================================================================
\aufgabe{B2: Erkenntnis -- Die Reihenfolge}

\vspace{0.2em}
\begin{center}
{\large $F_{\text{Haft}}$ \lsg{$>$} $F_{\text{Gleit}}$ \lsg{$>$} $F_{\text{Roll}}$}
\end{center}

\vspace{0.1em}
\begin{merkebox}[title=K2: Merksatz -- Die drei Reibungsarten]
\small
Die \lsg{Haftreibung} ist am größten, die \lsg{Rollreibung} am kleinsten.
$\rightarrow$ Deshalb ist das Erfinden des \lsg{Rades} so wichtig: Rollreibung statt Gleitreibung!
\end{merkebox}

\vspace{0.2em}

% ============================================================================
% B3/B4: EINFLUSSFAKTOREN
% ============================================================================
\aufgabe{B3/B4: Einflussfaktoren der Reibung}

\vspace{0.2em}
\begin{tabular}{ll}
B3: & Mehr Gewichtskraft = \lsg{mehr / größere} Reibung \\[0.3em]
B4: & Rauere Oberfläche = \lsg{mehr / größere} Reibung \\
\end{tabular}

\vspace{0.2em}

% ============================================================================
% K4: MESSWERTTABELLE + MERKSATZ EINFLUSSFAKTOREN
% ============================================================================
\begin{protokollbox}[title=K4: Messwerte -- Reibungskraft bei verschiedenen Gewichten]
\small
\renewcommand{\arraystretch}{1.3}
\begin{tabular}{|p{3.5cm}|p{3.5cm}|p{3.5cm}|}
\hline
\textbf{$F_\text{N}$ (Gewichtskraft)} & \textbf{$F_\text{R}$ (Reibungskraft)} & \textbf{$F_\text{R} \div F_\text{N}$} \\
\hline
5 N & \lsg{2,5} N & \lsg{0,50} \\
\hline
10 N & \lsg{5,0} N & \lsg{0,50} \\
\hline
15 N & \lsg{7,5} N & \lsg{0,50} \\
\hline
\end{tabular}

\vspace{0.2em}
\textbf{Erkenntnis:} Der Quotient ist immer \lsg{gleich (0,50)} $\rightarrow$ Das ist die \textbf{Reibungszahl} $\mu$!
\end{protokollbox}

\vspace{0.2em}

\begin{merkebox}[title=K4b: Merksatz -- Einflussfaktoren der Reibung]
\small
Die Reibungskraft hängt ab von: (1) der \lsg{Normalkraft (Gewichtskraft)} und (2) der \lsg{Oberflächenbeschaffenheit (Reibungszahl $\mu$)}.
\end{merkebox}

\vspace{0.2em}

% ============================================================================
% T2: ERWUENSCHT/UNERWUENSCHT
% ============================================================================
\begin{protokollbox}[title=T2: Erwünschte und unerwünschte Reibung]
\small
\renewcommand{\arraystretch}{1.3}
\begin{tabular}{|p{5.8cm}|p{5.8cm}|}
\hline
\textbf{Erwünscht (nützlich)} & \textbf{Unerwünscht (störend)} \\
\hline
1. \lsg{Bremsen} & 1. \lsg{Motor verschleißt} \\[0.3em]
2. \lsg{Schuhe rutschen nicht} & 2. \lsg{Energie geht verloren} \\[0.3em]
3. \lsg{Schrauben halten} & 3. \lsg{Gelenke nutzen ab} \\
\hline
\end{tabular}
\end{protokollbox}

\vspace{0.2em}

\aufgabe{A1: Fahrrad}

Erwünscht: \lsg{Bremsen (Bremsbeläge auf Felge/Scheibe), Reifen auf Straße}

Unerwünscht: \lsg{Kette, Lager, Tretlager (braucht Öl/Fett)}

\vspace{0.2em}
\begin{merkebox}[title=K5: Merksatz -- Reibung verringern]
\small
Reibung kann verringert werden durch: \lsg{Schmiermittel (Öl, Fett)}, \lsg{glatte Oberflächen} und \lsg{Kugellager}.
\end{merkebox}

\vspace{0.4em}

% ============================================================================
% K3: FORMEL + DREIECK
% ============================================================================
\begin{merkebox}[title=K3: Die Reibungsformel \stern{$\star\star$}]
\small
\textbf{Die Formel:} \quad {\large $F_\text{R} = \mu \cdot F_\text{N}$}

\vspace{0.3em}
\renewcommand{\arraystretch}{1.2}
\begin{tabular}{ll}
$F_\text{R}$ & = Reibungskraft in \lsg{N (Newton)} \\
$\mu$ & = Reibungszahl (``mü''), \lsg{keine} Einheit \\
$F_\text{N}$ & = Normalkraft in \lsg{N (Newton)} \\
\end{tabular}

\vspace{0.2em}
{\small\textbf{Normalkraft} = die Kraft, mit der ein Gegenstand auf die Unterlage drückt. Auf ebener Fläche: $F_\text{N} \approx F_\text{G}$.}

\vspace{0.3em}
\renewcommand{\arraystretch}{1.5}
\begin{tabular}{|p{2cm}|p{6cm}|}
\hline
\textbf{Gesucht} & \textbf{Formel} \\
\hline
$F_\text{R}$ & $F_\text{R} =$ \lsg{$\mu \cdot F_\text{N}$} \\
\hline
$\mu$ & $\mu =$ \lsg{$\dfrac{F_\text{R}}{F_\text{N}}$} \\
\hline
$F_\text{N}$ & $F_N =$ \lsg{$\dfrac{F_\text{R}}{\mu}$} \\
\hline
\end{tabular}
\end{merkebox}

\vspace{0.4em}

% ============================================================================
% RECHENAUFGABEN
% ============================================================================
\aufgabe{Rechenaufgaben -- Lösungen}

\vspace{0.3em}
\textbf{\stern{$\star$} R1:}

\lsg{Geg.: $\mu = 0{,}4$, $F_N = \SI{20}{N}$ \quad Ges.: $F_\text{R}$}

\lsg{Formel: $F_\text{R} = \mu \cdot F_\text{N}$}

\lsg{Einsetzen: $F_\text{R} = 0{,}4 \cdot \SI{20}{N} = \SI{8}{N}$}

\vspace{0.4em}

\textbf{\stern{$\star\star$} R2:}

\lsg{Geg.: $\mu = 0{,}8$, $F_N = \SI{5000}{N}$ \quad Ges.: $F_\text{R}$}

\lsg{Formel: $F_\text{R} = \mu \cdot F_\text{N}$}

\lsg{Einsetzen: $F_\text{R} = 0{,}8 \cdot \SI{5000}{N} = \SI{4000}{N}$}

\vspace{0.4em}

\textbf{\stern{$\star\star\star$} R3:}

\lsg{Geg.: $F_\text{R} = \SI{12}{N}$, $F_N = \SI{40}{N}$ \quad Ges.: $\mu$}

\lsg{Formel: $\mu = \dfrac{F_\text{R}}{F_\text{N}}$}

\lsg{Einsetzen: $\mu = \dfrac{\SI{12}{N}}{\SI{40}{N}} = 0{,}3$}

\vspace{0.4em}

% ============================================================================
% A2: BREMSWEG
% ============================================================================
\begin{protokollbox}[title=A2: Anwendung -- Bremsweg]
\small
a) Welche Reibungsart ist beim Bremsen wichtig? \lsg{Haftreibung}

\vspace{0.2em}
b) Auf welcher Straße ist der Bremsweg am längsten? \lsg{Vereiste Straße}

\vspace{0.2em}
c) Erkläre in einem Satz, warum:

\lsg{Je weniger Reibung zwischen Reifen und Straße, desto länger der Bremsweg. Auf Eis ist $\mu$ sehr klein.}

\vspace{0.2em}
\textbf{K6 Merksatz:} Je \lsg{geringer} die Reibung, desto \lsg{länger} der Bremsweg!
\end{protokollbox}

\end{document}
